% =====================================================================
% 3 Referencial teórico e revisão da literatura
% Fontes: paper/secoes/03-teoria.tex; paper/secoes/01-introducao.tex
% (antecedentes); docs/ars/13 (Calzada); docs/ars/16 (Denteh & Kédagni).
% =====================================================================

\chapter{Referencial teórico e revisão da literatura}
\label{cap:referencial}

\section{A deriva como externalidade negativa}
\label{sec:externalidade}

O agrotóxico que deriva sobre o vizinho é um custo imposto a terceiros que não
passa pelo sistema de preços. O produtor que contrata a aplicação aérea
internaliza o custo do serviço e do produto, e não o custo em saúde suportado
pela população a favor do vento. O resultado privado é aplicação acima do nível
socialmente eficiente.

As duas respostas canônicas divergem no instrumento. A pigouviana internaliza o
custo por tributo igual ao dano marginal. A coaseana observa que, com direitos
de propriedade bem definidos e custos de transação baixos, as partes negociam a
solução eficiente \cite{coase1960}, e desloca a pergunta para quem detém o
direito: o de pulverizar ou o de não ser pulverizado. A barganha coaseana não
se aplica a este caso. A deriva atinge população rural difusa, numerosa e não
organizada, com dano probabilístico, latente e de difícil atribuição
individual. Os custos de transação são proibitivos, e é isso que justifica a
intervenção regulatória e leva à discussão sobre o instrumento.

\section{Escolha de instrumento}
\label{sec:instrumento}

Sob incerteza a respeito dos custos de controle, um instrumento de preço e um
de quantidade não são equivalentes em bem-estar esperado \cite{weitzman1974}.
Na aproximação linear-quadrática, qual deles domina depende das inclinações
relativas do benefício marginal e do custo marginal do controle: se o benefício
marginal é mais inclinado, domina a quantidade; se o custo marginal é mais
inclinado, domina o preço.

Três qualificações da literatura importam para este caso, e o
Capítulo~\ref{cap:ensaio2} as desenvolve. A primeira é que proibir um método de
aplicação é um padrão tecnológico, e padrões sobre insumo ou tecnologia têm
efeitos distintos dos de padrões sobre emissão \cite{helfand1991}. A segunda é
que a deriva é poluição de fonte difusa: o regulador não observa a contribuição
de cada produtor, e o instrumento viável depende do que se consegue monitorar
\cite{segerson1988}. Sob heterogeneidade entre produtores, a taxa tende a ser
mais eficiente e a regulação por quantidade tende a gerar maior lucro ao setor
regulado \cite{helfand1995}. A terceira é que a literatura específica sobre
pesticidas trata dos custos de bem-estar e dos efeitos distributivos da
regulação \cite{lichtenberg1988,fernandezcornejo1998} e registra cenários em
que cotas superam taxas \cite{skevas2012}, além de simulação de banimento de
molécula na Europa \cite{bocker2020glyphosate}.

Há uma camada normativa que o cálculo de bem-estar esperado não capta.
Precificar uma externalidade de saúde supõe que o dano é comensurável com
dinheiro e que a compensação é resposta adequada. Proibi-la supõe um direito a
não ser exposto, anterior ao cálculo. \citet{hitzig2020normative} chama de
\emph{normative gap} a distância entre os objetivos de quem formula a política
e os objetivos que o economista embute no desenho do instrumento. A lei
cearense e a decisão do STF que a validou raciocinam em termos de saúde,
ambiente e exposição involuntária de terceiros. Ordenar proibição e taxa apenas
por bem-estar esperado responde a uma pergunta que o legislador não formulou.

\section{Evidência empírica sobre agrotóxicos e saúde ao nascer}
\label{sec:evidencia}

A evidência econômica se organiza em torno de duas perguntas: o que uma
substância faz à saúde e o que um método de aplicação faz à saúde.

Do lado da substância, \citet{dias2023down} são a referência brasileira.
Combinam o ganho potencial de produtividade trazido pela soja transgênica a
cada município com a direção do fluxo d'água na bacia e documentam piora dos
desfechos ao nascer a jusante. \citet{reynier2025glyphosate} estimam, para os
Estados Unidos, redução de 23 a 32 g no peso ao nascer à intensidade média de
exposição, concentrada no decil inferior de peso esperado.
\citet{larsen2017agricultural} encontram, no Vale de San Joaquin, efeitos
adversos de 5 a 9\% apenas entre os expostos ao topo da distribuição, o que
recomenda tratar a exposição como contínua e desconfiar de linearidade.
\citet{frank2024ecosystem} exploram um choque que eleva o uso de inseticida, a
perda de morcegos predadores, e encontram aumento da mortalidade infantil nos
condados atingidos.

Do lado do método, \citet{camacho2017health} avaliam as consequências
sanitárias da aspersão aérea de herbicida sobre cultivos ilícitos na Colômbia.
O análogo químico e de cultura mais próximo é \citet{calzada2023bananas}, sobre
a fumigação aérea de bananais no Equador, com déficit de peso ao nascer de 80 a
150 g sob exposição alta, segundo o resumo do artigo. Esse intervalo é a
referência de magnitude usada no exercício de calibração da
Seção~\ref{sec:diluicao}.

Nenhum desses trabalhos avalia uma proibição de método. É essa a lacuna que o
Ensaio 1 ocupa.

\section{Literatura metodológica}
\label{sec:metodo-lit}

A proibição entrou em vigor de uma só vez em todo o estado. Não há variação na
data de adoção, e os estimadores desenhados para adoção escalonada
\cite{callaway2021multiple,goodmanbacon2021,dechaisemartin2020} não se aplicam.
A variação disponível é de intensidade. \citet{callaway2024continuous} mostram
que diferenças em diferenças com tratamento contínuo não recupera a curva
dose-resposta sob as hipóteses habituais e distinguem dois parâmetros: o efeito
no nível de dose recebido, identificado sob tendências paralelas usuais, e a
curva causal, que exige tendências paralelas fortes.

Quatro outras peças sustentam o desenho. \citet{roth2023trending} mostra que um
teste de tendências anteriores pode passar por falta de poder.
\citet{rambachan2023credible} oferecem intervalos que admitem violações
limitadas do paralelismo. \citet{conley2011} e \citet{cameron2008wild} tratam
da inferência com poucos grupos, com as ressalvas de \citet{canay2021wild}. E
\citet{denteh2022misclassification} estudam diferenças em diferenças com
tratamento mal classificado: quando não há falso negativo, o parâmetro estimado
é o efeito verdadeiro multiplicado pelo valor preditivo positivo da
classificação. Esse resultado organiza a leitura dos resultados preliminares
(Seção~\ref{sec:diluicao}).
